%!TeX program = xelatex
\documentclass{SYSUReport}
\usepackage{tabularx} % 在导言区添加此行
\usepackage{float}
% 根据个人情况修改
\headl{}
\headc{}
\headr{并行程序设计与算法实验}
\lessonTitle{并行程序设计与算法实验}
\reportTitle{Lab7-MPI并行应用}
\stuname{xxx}
\stuid{xxx}
\inst{计算机学院}
\major{xxx}
\date{2025年5月7日}

\begin{document}

% =============================================
% Part 1: 封面
% =============================================
\cover
\thispagestyle{empty} % 首页不显示页码
\clearpage

% % =============================================
% % Part 4： 正文内容
% % =============================================
% % 重置页码，并使用阿拉伯数字
% \pagenumbering{arabic}
% \setcounter{page}{1}

%%可选择这里也放一个标题
%\begin{center}
%    \title{ \Huge \textbf{{标题}}}
%\end{center}

\section{实验目的}
\begin{itemize}
   \item 验证并行FFT的加速效果与可扩展性
    \item 评估数据打包对通信性能的优化作用
    \item 分析并行规模对内存消耗的影响
\end{itemize}

\section{实验内容}
\begin{itemize}
    \item \textbf{串行FFT分析与并行化改造：}
    \begin{itemize}
        \item 阅读并理解提供的串行傅里叶变换代码 (fft\_serial.cpp)。
        \item 使用MPI (Message Passing Interface) 对串行FFT代码进行并行化改造，可能需要对原有代码结构进行调整以适应并行计算模型。
    \end{itemize}
    \item \textbf{MPI数据通信优化：}
    \begin{itemize}
        \item 研究并应用MPI数据打包技术（例如使用 \texttt{MPI\_Pack}/\texttt{MPI\_Unpack} 或 \texttt{MPI\_Type\_create\_struct}）来对通信数据进行重组，以优化消息传递效率。
    \end{itemize}
    \item \textbf{程序性能与内存分析：}
    \begin{itemize}
        \item \textbf{性能分析：}
        \begin{itemize}
            \item 分析在不同并行规模（即不同的进程数量）以及不同问题规模（即输入数据N的大小）条件下，并行FFT程序的性能表现（例如，通过计算加速比和并行效率来衡量）。
            \item 分析数据打包技术对于并行程序整体性能的具体影响。
        \end{itemize}
        \item \textbf{内存消耗分析：}
        \begin{itemize}
            \item 使用Valgrind工具集中的Massif工具来采集并分析并行程序在不同配置下的内存消耗情况。
            \item 在Valgrind命令中增加 \texttt{--stacks=yes} 参数以采集程序运行时栈内内存的消耗情况。
            \item 利用 \texttt{ms\_print} 工具将Massif输出的日志 (massif.out.pid) 可视化或转换为可读格式，分析内存消耗随程序运行时间的变化，特别是关注峰值内存消耗。
        \end{itemize}
    \end{itemize}
\end{itemize}


\section{实验结果与分析}
\subsection{并行FFT的实现与正确性验证}
\begin{itemize}
    \item 描述你的并行FFT算法设计和实现的关键点。
    \item 回答：
\end{itemize}

\subsection{数据打包优化}
\begin{itemize}
    \item 描述你所采用的数据打包方法及其实现。
    \item 回答：
\end{itemize}

\subsection{性能分析}
\begin{itemize}
    \item \textbf{不同问题规模 (N) 和并行规模（进程数 P）下的运行时间：}
    \begin{itemize}
        \item 请以表格形式展示在不同 N 和 P 组合下的原始运行时间。请根据你的实际实验情况填写数据。
      \begin{table}[h!]
        \centering
        \caption{不同问题规模 (N) 和并行规模 (P) 下的运行时间 (单位：秒)}
        \label{tab:runtime_results_optimized}
        \begin{tabular}{@{}lccccc@{}}
        \toprule
        \multirow{2}{*}{问题规模 (N)} & \multicolumn{5}{c}{并行规模 (进程数 P)} \\
        \cmidrule(l){2-6}
                                  & P=1 (串行) & P=2    & P=4    & P=8    & P=16 \\ \midrule
        $N_1 = 2^{16} $   &            &        &        &        &             \\
        $N_2 = 2^{18}$  &            &        &        &        &             \\
        $N_3 = 2^{20}$&            &        &        &        &             \\
        % 如需更多N值或不同的P值，可在此处调整或添加行/列
        % 例如:
        % $N_4 = 2^{22}$&            &        &        &        &             \\
        \bottomrule
        \end{tabular}
        \end{table}
    \end{itemize}
    \item \textbf{加速比 (Speedup) 分析：}
    \begin{itemize}
        \item 根据运行时间计算不同配置下的加速比 $S_p = T_{serial} / T_{parallel}$。
        \begin{table}[h!]
        \centering
        \caption{不同问题规模 (N) 和并行规模 (P) 下的加速比 ($S_p$)}
        \label{tab:speedup_results}
        \begin{tabular}{@{}lccccc@{}}
        \toprule
        \multirow{2}{*}{问题规模 (N)} & \multicolumn{5}{c}{并行规模 (进程数 P)} \\
        \cmidrule(l){2-6}
                                  & P=1 (基准) & P=2    & P=4    & P=8    & P=16\\ \midrule
        $N_1 = 2^{16}$   & 1          &        &        &        &             \\
        $N_2 = 2^{18}$  & 1          &        &        &        &             \\
        $N_3 = 2^{20} $& 1          &        &        &        &             \\
        % 如在运行时间记录中使用了其他N值，请在此处保持一致
        % 例如:
        % $N_4 = 2^{22}$& 1          &        &        &        &             \\
        \bottomrule
        \end{tabular}
        \end{table}
        \item 分析加速比，讨论其是否符合预期，并解释原因（例如，通信开销、负载均衡等）。
    \end{itemize}
    \item \textbf{数据打包对性能的影响：}
    \begin{itemize}
        \item 分析数据打包带来的性能提升或可能引入的额外开销。
        \item 回答：
    \end{itemize}
\end{itemize}

\subsection{内存消耗分析 (Valgrind Massif)}
\begin{itemize}
    \item 展示并分析由 \texttt{ms\_print} 生成的内存消耗图表或关键数据点。
    \item 你的图表：
    \item 分析不同并行规模（进程数）对程序峰值内存消耗、总内存分配量的影响。
    \item 回答：
\end{itemize}


注：实验报告格式参考本模板，可在此基础上进行修改；实验代码以zip格式另提交；最终提交内容包括实验报告(pdf格式)和实验代码(zip压缩包格式)
\end{document}